% LaTeXML conversion preamble for PoRA.
% Replaces the print preamble (tufte-book / tcolorbox / minted / biblatex /
% pgfplots), none of which convert cleanly. Loads only LaTeXML-supported
% packages and re-defines the book's custom environments/commands using classic
% LaTeX2e primitives (NOT xparse: in TeX Live 2025 the xparse/hyperref LaTeXML
% bindings pull in the real expl3-code.tex, which LaTeXML 0.8.8 cannot parse).

\usepackage{amsmath}
\usepackage{amsthm}
\usepackage{amsfonts}
\usepackage{amssymb}
\usepackage{mathtools}
\usepackage{bm}
\usepackage{graphicx}
\usepackage{subcaption}
\usepackage{xcolor}
\usepackage{enumitem}
\usepackage{booktabs}
\usepackage{array}
\usepackage{xspace}
\usepackage{url}
\usepackage{float}

% bbm/accents fallbacks
\providecommand{\mathbbm}[1]{\mathbb{#1}}
\providecommand{\accentset}[2]{\overset{#1}{#2}}
\providecommand{\nicefrac}[2]{#1/#2}
% Without hyperref, LaTeXML doesn't linkify \href/\url — the converter linkifies
% URLs in post-processing. Emit the href target (as \url) so it gets picked up.
\providecommand{\href}[2]{#2 (\url{#1})}
\providecommand{\texorpdfstring}[2]{#1}

% ---- theorem-like environments (amsthm; good LaTeXML support) ----
\theoremstyle{plain}
\newtheorem{theorem}{Theorem}[section]
\theoremstyle{definition}
\newtheorem{example}{Example}[section]
\newtheorem{definition}[theorem]{Definition}
\newtheorem{codex}[theorem]{Code Exercise}

% ---- boxed environments -> pass-through blocks (styled by the site theme) ----
\newenvironment{hint}{\par\noindent\textit{Hint:}\enspace\ignorespaces}{\par}
\newenvironment{tcolorbox}[1][]{\par\noindent}{\par}
\newenvironment{definitionbox}[1][]{\par\noindent}{\par}
\newcommand{\adjustbox}[2]{#2}
% \resizebox/\scalebox around a (pre-rendered) tikz figure make LaTeXML emit a
% CSS scale() transform whose scaled height ISN'T reserved in layout, so the
% figure overflows onto its caption. Drop the box; the SVG then sizes to the
% column via CSS (max-width:100%).
\renewcommand{\resizebox}[3]{#3}
\newcommand{\scalebox}[2]{#2}

% ---- margin material (tufte) -> LaTeXML margin notes ----
\makeatletter
% \sidenote swallows up to two optional args, keeps the mandatory text.
\newcommand{\sidenote}{\@ifnextchar[{\sn@a}{\sn@a[]}}
\def\sn@a[#1]{\@ifnextchar[{\sn@b[#1]}{\sn@b[#1][]}}
\def\sn@b[#1][#2]#3{\marginpar{#3}}
\makeatother
\newcommand{\mysidenote}[1]{\marginpar{#1}}
\newcommand{\marginnote}[2][]{\marginpar{#2}}

% ---- citations ------------------------------------------------------------
% The print book renders \cite as a Tufte sidenote holding the *full* reference
% (biblatex \sidenote{\fullcite{...}}). We can't run biblatex under LaTeXML, so
% each citation emits a plain-text marker carrying its key list; scripts/bib.mjs
% (called from convert.mjs) parses references.bib and rewrites the markers into
% numbered margin sidenotes (\cite/\citep), inline "Author (Year)" + sidenote
% (\citet/\textcite), or an inline full reference (\fullcite). Markers are
% emitted as inline text (NOT \marginpar) so citations that appear inside a
% \sidenote/\footnote don't produce invalid nested margin notes.
\makeatletter
% \lx@cite{<handler>} discards up to two biblatex optional args
% ([prenote][postnote]) then applies <handler> to the mandatory {key list}.
\newcommand{\lx@cite}[1]{\@ifnextchar[{\lx@cite@i{#1}}{#1}}
\def\lx@cite@i#1[#2]{\@ifnextchar[{\lx@cite@ii{#1}}{#1}}
\def\lx@cite@ii#1[#2]{#1}
\newcommand{\lx@citemark}[1]{lxCITEP #1 lxCITEEND}
\newcommand{\lx@textmark}[1]{lxCITET #1 lxCITEEND}
\newcommand{\lx@fullmark}[1]{lxCITEF #1 lxCITEEND}
\renewcommand{\cite}{\lx@cite{\lx@citemark}}
\newcommand{\citep}{\lx@cite{\lx@citemark}}
\newcommand{\citet}{\lx@cite{\lx@textmark}}
\newcommand{\textcite}{\lx@cite{\lx@textmark}}
\newcommand{\fullcite}{\lx@cite{\lx@fullmark}}
\makeatother

\newcommand{\cref}[1]{\ref{#1}}
\newcommand{\Cref}[1]{\ref{#1}}
\newenvironment{marginfigure}[1][]{\begin{figure}}{\end{figure}}

% ---- csquotes replacements (avoids expl3) ----
\newcommand{\enquote}[1]{``#1''}
\newenvironment{displayquote}{\begin{quote}}{\end{quote}}

% ---- colored-table + stray tikz/pgf commands -> no-ops ----
\newcommand{\columncolor}[2][]{}
\newcommand{\cellcolor}[2][]{}
\newcommand{\rowcolor}[2][]{}
\providecommand{\usetikzlibrary}[1]{}
\providecommand{\usepgfplotslibrary}[1]{}
\providecommand{\pgfplotsset}[1]{}

% ---- misc print-only commands -> safe web equivalents ----
\newcommand{\dg}[1]{}
\newcommand{\colorcode}[1]{\texttt{#1}}
\newcommand{\writtenQ}{}
\newcommand{\codeQ}{}
\newcommand{\notebook}[2]{\texttt{#1}}
\newcommand{\codeword}[1]{\texttt{#1}}
\newcommand{\nlt}{\par\noindent}
\providecommand{\blankpage}{}
\providecommand{\notessection}[1]{\section*{#1}}

% ---- code listings -> lstlisting (syntax classes coloured in latexml.css) ----
% The print book uses minted; here we route every code environment through
% `listings` with a language set, so LaTeXML emits ltx_lst_keyword /
% ltx_lst_string / ltx_lst_comment classes (styled to match the book's
% deep-blue keywords / green strings). columns=fullflexible + keepspaces keeps
% indentation; escapeinside={||} makes the book's |$...$| render as inline math.
\usepackage{listings}
\lstset{basicstyle=\ttfamily,columns=fullflexible,keepspaces=true,%
  showstringspaces=false,breaklines=true,otherkeywords={self}}
% minted: optional [<lstopts>] then mandatory {<language>}
% (e.g. [escapeinside=||]{python}, or {bash}).
\lstnewenvironment{minted}[2][]{\lstset{language=#2,#1}}{}
\lstnewenvironment{python}[1][]{\lstset{language=Python,escapeinside={||},#1}}{}
\lstnewenvironment{pythonnoborder}[1][]{\lstset{language=Python,escapeinside={||},#1}}{}
\lstnewenvironment{gencode}[1][]{\lstset{#1}}{}
\newcommand{\pythonexternal}[2][]{\lstinputlisting[language=Python,#1]{#2}}
\newcommand{\pythoninline}[1]{\lstinline!#1!}
% "listing" float used by the book for code blocks (print calls it "Algorithm").
% Numbered continuously per chapter (the converter sets \thelisting = "N.M");
% dropping the [section] reset avoids the counter restarting and colliding.
\newfloat{listing}{tbp}{lol}
\floatname{listing}{Algorithm}

% ---- algorithms: the book mixes algorithm2e and algpseudocode ----
\usepackage[ruled,vlined]{algorithm2e}

% Pre-rendered tikz figures are injected as a text token that the converter
% replaces with an <img> in post-processing (LaTeXML can't \includegraphics SVG).
\newcommand{\lxtikz}[1]{TKZFIG#1ENDTKZ}

% The book authored chapter headings one level down (the print preamble does
% \let\subsection\section). Replicate so \subsection becomes a top-level heading.
\let\subsubsection\subsection
\let\subsection\section

% shared math macros
% PoRA math macros, consolidated from the book's preamble.tex + utils/symbols.tex
% for LaTeXML conversion. Print-only macros (colors, tikz styles, listings,
% title formatting, citation-as-sidenote) are intentionally omitted; those are
% handled in latexml-preamble.tex. The two malformed \maximize[1]/\minimize[1]
% definitions in the original symbols.tex are fixed here.

% ---- operators ----
\DeclareMathOperator*{\argmax}{arg\,max}
\DeclareMathOperator*{\argmin}{arg\,min}
\DeclareMathOperator{\Exp}{Exp}

% ---- paired delimiters (mathtools) ----
\DeclarePairedDelimiter{\norm}{\lVert}{\rVert}
\DeclarePairedDelimiter{\rbr}{(}{)}

% ---- misc math shorthands (from preamble.tex) ----
\newcommand{\defn}{\coloneqq}
\newcommand{\inv}[1]{#1^{-1}}
\newcommand{\pd}[3][]{\frac{\partial^{#1}#2}{{\partial#3}^{#1}}}
\newcommand{\mytilde}{\raise.17ex\hbox{$\scriptstyle\mathtt{\sim}$}}
\newcommand{\tran}[1]{#1^\mathsf{T}}
\newcommand{\bmx}[2][1]{\begin{bmatrix}#2\end{bmatrix}}

% ---- bold/vector symbols ----
\newcommand{\bxi}{\bm{\xi}}
\renewcommand{\a}{\bm{a}}
\renewcommand{\v}{\bm{v}}
\newcommand{\e}{\bm{e}}
\newcommand{\x}{\bm{x}}
\newcommand{\bX}{\bm{X}}
\newcommand{\y}{\bm{y}}
\newcommand{\ac}{\bm{u}}
\newcommand{\z}{\bm{z}}
\newcommand{\p}{\bm{p}}
\newcommand{\blam}{\bm{\lambda}}
\newcommand{\bnu}{\bm{\nu}}
\newcommand{\bu}{\bm{u}}
\newcommand{\bg}{\bm{g}}
\newcommand{\q}{\bm{q}}
\newcommand{\bv}{\bm{v}}
\newcommand{\bmu}{\bm{\mu}}
\newcommand{\bSigma}{\bm{\Sigma}}
\newcommand{\m}{\bm{m}}
\newcommand{\w}{\bm{w}}
\newcommand{\bc}{\bm{c}}
\newcommand{\btheta}{\bm{\theta}}
\newcommand{\bphi}{\bm{\phi}}
\newcommand{\btau}{\bm{\tau}}
\newcommand{\f}{f}
\newcommand{\bd}{\bm{d}}

% ---- sets / blackboard / calligraphic ----
\newcommand{\R}{\mathbb{R}}
\newcommand{\C}{\mathcal{C}}
\newcommand{\graph}{\mathcal{G}}
\newcommand{\edge}{\mathcal{E}}
\newcommand{\node}{\mathcal{V}}
\newcommand{\startnode}{q_{\mathrm{S}}}
\newcommand{\goalnode}{q_{\mathrm{G}}}
\newcommand{\Space}{\mathcal{S}}
\newcommand{\E}{\mathbb{E}}
\newcommand{\X}{\mathcal{X}}
\newcommand{\U}{\mathcal{U}}
\newcommand{\dagname}{\text{DA\footnotesize{GGER}}}
\newcommand{\stateNoise}{Q}
\newcommand{\measNoise}{R}

% ---- from utils/symbols.tex ----
\newcommand{\bel}{\text{bel}}
\newcommand{\belpred}{\overline{\bel}}
\newcommand{\code}[1]{\texttt{#1}}
\newcommand{\controldim}{m}
\newcommand{\dynmodel}{f}
\newcommand{\dynJac}{F}
\renewcommand{\d}{d}
\newcommand{\definedas}{\coloneqq}
\newcommand{\expected}[2]{\mathbb{E}_{#1}\left[#2 \right]}
\newcommand{\genCa}{\bxi}
\newcommand{\given}{\mid}
\newcommand{\hamiltonian}{H}
\newcommand{\controlspace}{\mathcal{U}}
\newcommand{\lagrangian}{L}
\newcommand{\measmodel}{h}
\newcommand{\measJac}{H}
% fixed (originally malformed \maximize[1]/\minimize[1]):
\newcommand{\maximize}[1]{\underset{#1}{\text{maximize}}\quad}
\newcommand{\minimize}[1]{\underset{#1}{\text{minimize}}\:\quad}
\newcommand{\Oc}{\mathcal{O}}
\newcommand{\outputdim}{p}
\newcommand{\particleset}{\mathcal{P}}
\newcommand{\posreals}{\mathbb{R}_{\geq 0}}
\newcommand{\pluseq}{\mathrel{+}=}
\newcommand{\reals}{\mathbb{R}}
\newcommand{\statespace}{\mathcal{X}}
\newcommand{\statedim}{n}
\newcommand{\setof}[1]{\left\{ #1\right\}}
\newcommand{\subjectto}{\text{subject to} \quad}
\newcommand{\tup}[1]{\left(#1\right)}
\renewcommand{\u}{\bm{u}}
\newcommand{\vCol}[1]{\left[{#1}\right]^\top}

